% ===== PRÉAMBULE CANONIQUE — à coller VERBATIM en tête de CHAQUE .tex =====
\documentclass[11pt,a4paper]{article}
\usepackage[utf8]{inputenc}\usepackage[T1]{fontenc}\usepackage{lmodern}
\usepackage[french]{babel}
\usepackage{amsmath,amssymb,mathtools}\usepackage{icomma}
\usepackage{siunitx}\sisetup{locale=FR}  % \qty{}{}, \si{}, \ang{}
\usepackage{xcolor}\usepackage{colortbl}
\usepackage{tikz}\usepackage{pgfplots}\pgfplotsset{compat=1.18}
\usepackage[siunitx,europeanresistors]{circuitikz} % circuits (élec.)
\usepackage[most]{tcolorbox}\usepackage{enumitem}\usepackage{array,booktabs}
\usepackage[a4paper,margin=2.2cm]{geometry}
\usetikzlibrary{arrows.meta,calc,angles,quotes,positioning,patterns,%
  backgrounds,shapes.geometric,decorations.pathmorphing,decorations.markings,intersections}
\definecolor{primary}{RGB}{222,49,99}\definecolor{primarydark}{RGB}{150,22,55}
\definecolor{defcol}{RGB}{0,114,178}\definecolor{methcol}{RGB}{52,92,148}
\definecolor{excol}{RGB}{0,135,90}\definecolor{attncol}{RGB}{200,40,55}
\tcbset{boxbase/.style={enhanced,breakable,boxrule=0.9pt,arc=2.5pt,
  left=3mm,right=3mm,top=2.3mm,bottom=2.3mm,fonttitle=\bfseries,coltitle=white}}
% --- boîtes TUTORIEL (noms EXACTS, ne pas renommer) ---
\newtcolorbox{cle}[1][]{boxbase,colback=defcol!6,colframe=defcol,
  title={\textsc{À retenir}\ifx\relax#1\relax\relax\else\textmd{ : \itshape #1}\fi}}
\newtcolorbox{methode}[1][]{boxbase,colback=methcol!7,colframe=methcol,sharp corners,
  title={\textsc{Méthode}\ifx\relax#1\relax\relax\else\textmd{ --- \itshape #1}\fi}}
\newtcolorbox{exemplebox}[1][]{boxbase,colback=excol!5,colframe=excol,
  title={\textsc{Exemple}\ifx\relax#1\relax\relax\else\textmd{ : \itshape #1}\fi}}
\newtcolorbox{piege}[1][]{boxbase,colback=attncol!6,colframe=attncol,
  title={\textsc{Piège}\ifx\relax#1\relax\relax\else\textmd{ : \itshape #1}\fi}}
\newtcolorbox{remarque}[1][]{boxbase,colback=black!4,colframe=black!45,
  title={\textsc{Remarque}\ifx\relax#1\relax\relax\else\textmd{ : \itshape #1}\fi}}
% --- boîtes EXERCICE / CORRIGÉ (noms EXACTS, auto-comptées) ---
\newtcolorbox[auto counter]{exo}[1][]{boxbase,colback=primary!4,colframe=primary,
  title={\textsc{Exercice}~\thetcbcounter\ifx\relax#1\relax\relax\else\textmd{ --- \itshape #1}\fi}}
\newtcolorbox[auto counter]{corrige}[1][]{boxbase,colback=excol!4,colframe=excol,
  title={\textsc{Corrigé}~\thetcbcounter\ifx\relax#1\relax\relax\else\textmd{ --- \itshape #1}\fi}}
\newcommand{\vv}[1]{\vec{#1}}\newcommand{\ds}{\displaystyle}
\newcommand{\R}{\mathbb{R}}\newcommand{\N}{\mathbb{N}}\newcommand{\Z}{\mathbb{Z}}
% (un \usepackage{fancyhdr}+en-tête est possible ; il est ignoré côté web)
% =========================================================================
\begin{document}
\begin{corrige}[force avec un angle quelconque]
\[ F=B\,I\,L\,\sin\theta=0{,}2\times5\times0{,}20\times\sin\ang{30}
    =0{,}2\times5\times0{,}20\times0{,}5=\qty{0.1}{\newton}. \]
\end{corrige}

\begin{corrige}[retrouver le champ magnétique]
On isole $B$ dans $F=BIL$ (car $\theta=\ang{90}$), avec $L=\qty{0.15}{\metre}$ :
\[ B=\frac{F}{I\,L}=\frac{0{,}6}{8\times0{,}15}=\frac{0{,}6}{1{,}2}=\qty{0.5}{\tesla}. \]
\end{corrige}

\begin{corrige}[force entre deux fils parallèles]
\begin{enumerate}[label=\alph*)]
  \item $\dfrac{F}{L}=\dfrac{\mu_0 I_1 I_2}{2\pi d}=\dfrac{2\cdot10^{-7}\times5\times5}{0{,}02}
        =\dfrac{5\cdot10^{-6}}{0{,}02}=\qty{2.5e-4}{\newton\per\metre}.$
  \item Les courants sont de \textbf{même sens} : les fils \textbf{s'attirent}.
\end{enumerate}
\end{corrige}

\begin{corrige}[sens de la force avec $\vv{B}$ entrant]
Règle des trois doigts : \textbf{pouce} vers la droite ($I$), \textbf{index} vers l'arrière de la
feuille ($\vv{B}$ entrant), le \textbf{majeur} pointe alors \textbf{vers le haut}. La force de
Laplace est dirigée \textbf{vers le haut}.
\end{corrige}

\begin{corrige}[le rail de Laplace]
\begin{enumerate}[label=\alph*)]
  \item La barre est perpendiculaire à $\vv{B}$ ($\theta=\ang{90}$) :
        \[ F=B\,I\,L=0{,}3\times4\times0{,}08=\qty{9.6e-2}{\newton}\approx\qty{0.096}{\newton}. \]
  \item Trois doigts : \textbf{pouce} vers le haut ($I$ dans la barre), \textbf{index} vers l'arrière
        ($\vv{B}$ entrant), \textbf{majeur} vers la \textbf{gauche}. La barre est poussée et se déplace
        \textbf{vers la gauche} le long des rails. (Inverser $I$ ou $\vv{B}$ inverserait le sens : c'est
        le principe du haut-parleur et du moteur linéaire.)
\end{enumerate}
\end{corrige}

\end{document}
